\documentclass[12pt]{article}
\usepackage{amsmath, amssymb, amsthm}

\newtheorem{problem}{Problem}

\begin{document}

\section*{Probability}

\begin{problem}
Three closed boxes lie on a table. One box (you don’t know which) contains a \$1000 bill.
The others are empty. After paying an entry fee, you play the following game with the owner
of the boxes: you point to a box but do not open it; the owner then opens one of the two
remaining boxes and shows you that it is empty; you may now open either the box you first
pointed to or else the other unopened box, but not both. If you find the \$1000, you get to
keep it. Does it make any difference which box you choose? What is a fair entry fee for this
game?
\end{problem}

\begin{problem}
You are dealt two cards face down from a shuffled deck of 8 cards consisting of the four
queens and four kings from a standard bridge deck. The dealer looks at both of your two
cards (without showing them to you) and tells you (truthfully) that at least one card is a
queen. What is the probability that you have been given two queens? What is this probability
if the dealer tells you instead that at least one card is a red queen? What is this probability if
the dealer tells you instead that at least one card (or exactly one card) is the queen of hearts?
\end{problem}

\begin{problem}
An unfair coin (probability $p$ of showing heads) is tossed $n$ times. What is the probability
that the number of heads will be even?
\end{problem}

\begin{problem}
Two persons agreed to meet in a definite place between noon and one o’clock. If either person
arrives while the other is not present, he or she will wait for up to 15 minutes. Calculate the
probability that the meeting will occur, assuming that the arrival times are independent and
uniformly distributed between noon and one o’clock.
\end{problem}

\begin{problem}
Real numbers are chosen at random from the interval $[0,1]$. If after choosing the $n$th number
the sum of the numbers so chosen first exceeds 1, show that the expected or average value
for $n$ is $e$.
\end{problem}

\begin{problem}
Let $\alpha$ and $\beta$ be given positive real numbers with $\alpha < \beta$. If two points are selected at random
from a straight line segment of length $\beta$, what is the probability that the distance between
them is at least $\alpha$?
\end{problem}

\begin{problem}
Two real numbers $x$ and $y$ are chosen at random in the interval $(0,1)$ with respect to the
uniform distribution. What is the probability that the closest integer to $x/y$ is even? Express
the answer in the form $r + s\pi$, where $r$ and $s$ are rational numbers.
\end{problem}

\begin{problem}
Four points are chosen uniformly and independently at random in the interior of a given
circle. Find the probability that they are the vertices of a convex quadrilateral.
\end{problem}

\begin{problem}
Let $k$ be a positive integer. Suppose that the integers $1,2,3,\dots,3k+1$ are written down in
random order. What is the probability that at no time during this process, the sum of the
integers that have been written up to that time is a positive integer divisible by 3? Your
answer should be in closed form, but may include factorials.
\end{problem}

\begin{problem}
Let $(x_1, x_2, \dots, x_n)$ be a point chosen at random from the $n$-dimensional region defined by
$0 < x_1 < x_2 < \cdots < x_n < 1$. Let $f$ be a continuous function on $[0,1]$ with $f(1)=0$. Set
$x_0 = 0$ and $x_{n+1} = 1$. Show that the expected value of the Riemann sum
\[
\sum_{i=0}^n (x_{i+1}-x_i) f(x_{i+1})
\]
is
\[
\int_0^1 f(t)P(t)\, dt,
\]
where $P$ is a polynomial of degree $n$, independent of $f$, with $0 \leq P(t) \leq 1$ for $0 \leq t \leq 1$.
\end{problem}

\begin{problem}
Choose $n$ points $x_1, \dots, x_n$ at random from the unit interval $[0,1]$. Let $p_n$ be the probability
that $x_i+x_{i+1} \leq 1$ for all $1 \leq i \leq n-1$. Find a simple expression for
\[
\sum_{n \geq 0} p_n x^n = 1 + x + \tfrac{1}{2}x^2 + \tfrac{1}{3}x^3 + \cdots.
\]
\end{problem}

\begin{problem}
A dart, thrown at random, hits a square target. Assuming any two parts of the target of
equal area are equally likely to be hit, find the probability that the point hit is nearer to the
center than to any edge. Express your answer in the form
\[
\frac{a\sqrt{b}+c}{d},
\]
where $a,b,c,d$ are integers.
\end{problem}

\begin{problem}
If $\alpha$ is an irrational number, $0<\alpha<1$, is there a finite game with an honest coin such
that the probability of one player winning the game is $\alpha$? (An honest coin has probability
$1/2$ for both heads and tails. A game is finite if, with probability 1, it must end in a finite
number of moves.)
\end{problem}

\begin{problem}
Let $C$ be the unit circle $x^2+y^2=1$. A point $p$ is chosen randomly on the circumference
$C$ and another point $q$ is chosen randomly from the interior of $C$, independently and uniformly.
Let $R$ be the rectangle with sides parallel to the axes with diagonal $pq$. What is the probability
that no point of $R$ lies outside of $C$?
\end{problem}

\begin{problem}
Let $p_n$ be the probability that $c+d$ is a perfect square when the integers $c$ and $d$ are selected
independently at random from the set $\{1,2,\dots,n\}$. Show that
\[
\lim_{n\to\infty} \big( p_n \sqrt{n} \big)
\]
exists, and express this limit in the form $r(\sqrt{s}-t)$, where $s,t$ are integers and $r$ is a rational number.
\end{problem}

\begin{problem}
The points $1,2,\dots,1000$ are paired up at random to form 500 intervals $[i,j]$. What is the
probability that among these intervals is one which intersects all the others?
\end{problem}

\begin{problem}
The temperatures in Chicago and Detroit are $x^\circ$ and $y^\circ$, respectively. We are given:
\begin{enumerate}
\item $P(x^\circ = 70^\circ)$,
\item $P(y^\circ = 70^\circ)$,
\item $P(\max(x^\circ,y^\circ)=70^\circ)$.
\end{enumerate}
Determine $P(\min(x^\circ,y^\circ)=70^\circ)$.
\end{problem}

\begin{problem}
In the Massachusetts MEGABUCKS lottery, six distinct integers from 1 to 36 are selected
each week, uniformly at random. If $N_{\max}$ denotes the largest of the six numbers, find the
expected value of $N_{\max}$.
\end{problem}

\begin{problem}
\begin{enumerate}
\item[(a)] A fair die is tossed repeatedly. Let $p_n$ be the probability that after some number of
tosses the sum of the numbers that have appeared is $n$. (For instance, $p_1 = 1/6$ and
$p_2 = 7/36$.) Find $\lim_{n\to\infty} p_n$.

\item[(b)] More generally, suppose that a “die” has infinitely many faces, marked $1,2,\dots$. When
the die is thrown, the probability is $a_i$ that face $i$ appears (with $\sum_{i=1}^\infty a_i=1$). 
Let $p_n$ be as in (a), and find $\lim_{n\to\infty} p_n$. Assume that there does not exist 
$k>1$ such that if $a_i \neq 0$, then $k \mid i$.
\end{enumerate}
\end{problem}

\begin{problem}
Suppose that each of $n$ people write down the numbers $1,2,3$ in random order in one column
of a $3\times n$ matrix, independently and uniformly. Let the row sums $a,b,c$ be rearranged (if
necessary) so that $a \leq b \leq c$. Show that for some $n \geq 1995$, it is at least four times
as likely that both $b=a+1$ and $c=a+2$ as that $a=b=c$.
\end{problem}

\begin{problem}
At time $t=1$ choose two numbers $x_1,y_1$ uniformly and independently from $[0,1]$. At time
$t=2$ choose two further numbers $x_2,y_2$, etc. What is the expected time $n$ at which
\[
\sum_{i=1}^n (x_i^2 + y_i^2) > 1
\]
for the first time?
\end{problem}

\begin{problem}
A fair coin is flipped until the number of heads exceeds the number of tails. What is the
expected number of flips?
\end{problem}

\begin{problem}
An integer $n$, unknown to you, has been randomly chosen in the interval $[1,2002]$ with uniform
probability. Your objective is to select $n$ in an odd number of guesses. After each incorrect
guess, you are informed whether $n$ is higher or lower, and you must guess again among the
remaining feasible numbers. Show that you have a strategy such that the chance of winning
is greater than $2/3$.
\end{problem}

\begin{problem}
A deck of cards (26 red and 26 black) is shuffled, and the cards are turned face up one at a
time. At any point before the last card, you can say ``stop.'' If the next card is red, you win \$1;
if it is black, you win nothing. What is your best strategy? In particular, is there a strategy
with expected value greater than 50 cents?
\end{problem}

\begin{problem}
In the previous problem, suppose that you start with \$1 and after each card is shown you can
bet (at even odds) on red or black, any fraction of your current wealth. You can guarantee
that you end up with \$2 by waiting until one card remains before you bet. Can you guarantee
that you will end up with more than \$2? If so, what is the maximum amount you can be sure
of winning?
\end{problem}

\begin{problem}
Alice takes two slips of paper and writes a different integer on each. Bob chooses one slip
and looks at it. He may keep it or exchange it for the other slip. If he ends up with the larger
integer, he wins. Is there a strategy for Bob which gives him a probability of more than 50\%
of winning?
\end{problem}

\begin{problem}
Suppose in the previous problem that two real numbers in $[0,1]$ are chosen uniformly at random.
Alice looks at the two numbers and then decides which one to show Bob. If Alice chooses
optimally, can Bob do better than break even? What are the optimal strategies of Bob and Alice?
\end{problem}

\begin{problem}
Values $a_1, \dots, a_{2013}$ are chosen independently and at random from the set $\{1,\dots,2013\}$. What
is the expected number of distinct values in $\{a_1,\dots,a_{2013}\}$?
\end{problem}

\begin{problem}
Let $a,b,c,d,e,f$ be integers selected from $\{1,2,\dots,100\}$, uniformly and with replacement. Set
\[
M = a + 2b + 4c + 8d + 16e + 32f.
\]
What is the expected value of the remainder when $M$ is divided by 64?
\end{problem}

\begin{problem}
The number 2019 is written on a blackboard. Every minute, if the number $a$ is on the board,
Evan erases it and replaces it with a number chosen uniformly at random from
\[
\{0,1,2,\dots,\lceil 2.01a \rceil\}.
\]
Is there an integer $N$ such that the board reads 0 after $N$ steps with probability at least 99\%?
\end{problem}

\begin{problem}
Let $A$ be a $\sigma$-algebra with countable cardinality. Prove that $A$ is finite and determine the
possible values of its cardinality.
\end{problem}

\begin{problem}
Kevin has $2^n-1$ cookies, each labeled with a unique nonempty subset of $\{1,2,\dots,n\}$. Each
day, he chooses one cookie uniformly at random from those not yet eaten. Then he eats that
cookie and all remaining cookies labeled with a subset of that cookie. Determine the expected
number of days until all cookies are gone.
\end{problem}

\section*{Hidden Independence and Uniformity}

\begin{problem}
Slips of paper with the numbers from 1 to 99 are placed in a hat. Five numbers are randomly drawn out of the hat one at a time (without replacement). What is the probability that the numbers are chosen in increasing order?
\end{problem}

\begin{problem}
In how many ways can a positive integer $n$ be written as a sum of positive integers, taking order into account? For instance, $4$ can be written as a sum in the eight ways: 
\[
4 = 3 + 1 = 1 + 3 = 2 + 2 = 2 + 1 + 1 = 1 + 2 + 1 = 1 + 1 + 2 = 1 + 1 + 1 + 1.
\]
\end{problem}

\begin{problem}
\textbf{\textcolor{red}{\heart}} How many $8 \times 8$ matrices of 0’s and 1’s are there, such that every row and column contains an odd number of 1’s?
\end{problem}

\begin{problem}
Let $f(n)$ be the number of ways to take an $n$-element set $S$, and, if $S$ has more than one element, to partition $S$ into two disjoint nonempty subsets $S_1$ and $S_2$, then to take one of the sets $S_1$, $S_2$ with more than one element and partition it into two disjoint nonempty subsets $S_3$ and $S_4$, then to take one of the sets with more than one element not yet partitioned and partition it into two disjoint nonempty subsets, etc., always taking a set with more than one element that is not yet partitioned and partitioning it into two nonempty disjoint subsets, until only one-element subsets remain. 

For example, we could start with $12345678$ (short for $\{1,2,3,4,5,6,7,8\}$), then partition it into $126$ and $34578$, then partition $34578$ into $4$ and $3578$, then $126$ into $6$ and $12$, then $3578$ into $37$ and $58$, then $58$ into $5$ and $8$, then $12$ into $1$ and $2$, and finally $37$ into $3$ and $7$. (The order we partition the sets is important; for instance, partitioning $1234$ into $12$ and $34$, then $12$ into $1$ and $2$, and then $34$ into $3$ and $4$, is different from partitioning $1234$ into $12$ and $34$, then $34$ into $3$ and $4$, and then $12$ into $1$ and $2$. However, partitioning $1234$ into $12$ and $34$ is the same as partitioning it into $34$ and $12$.)

Find a simple formula for $f(n)$. For instance, $f(1) = 1$, $f(2) = 1$, $f(3) = 3$, and $f(4) = 18$.
\end{problem}

\begin{problem}
Fix positive integers $n$ and $k$. Find the number of $k$-tuples $(S_1,S_2,\ldots,S_k)$ of subsets $S_i$ of $\{1,2,\ldots,n\}$ subject to each of the following conditions:
    \begin{enumerate}
        \item $S_1 \subseteq S_2 \subseteq \cdots \subseteq S_k$
        \item The $S_i$’s are pairwise disjoint.
        \item $S_1 \cap S_2 \cap \cdots \cap S_k = \varnothing$
        \item $S_1 \subseteq S_2 \supseteq S_3 \subseteq S_4 \supseteq S_5 \subseteq \cdots S_k$\\
        (The symbols $\subseteq$ and $\supseteq$ alternate.)
    \end{enumerate}
\end{problem}

\begin{problem}
Let $p$ be a prime number and $1 \leq k \leq p-1$. How many $k$-element subsets $\{a_1,\ldots,a_k\}$ of $\{1,2,\ldots,p\}$ are there such that $a_1+\cdots+a_k \equiv 0 \pmod{p}$?
\end{problem}

\begin{problem}
\textbf{\textcolor{red}{\heart}} Let $\pi$ be a random permutation of $1,2,\ldots,n$. Fix a positive integer $1 \leq k \leq n$. What is the probability that in the disjoint cycle decomposition of $\pi$, the length of the cycle containing $1$ is $k$? In other words, what is the probability that $k$ is the least positive integer for which $\pi^k(1)=1$?
\end{problem}

\begin{problem}
Let $\pi$ be a random permutation of $1,2,\ldots,n$. What is the probability that $1$ and $2$ are in the same cycle of $\pi$?
\end{problem}

\begin{problem}
Choose $n$ real numbers $x_1,\ldots,x_n$ uniformly and independently from the interval $[0,1]$. What is the expected value of $\min_i x_i$, the minimum of $x_1,\ldots,x_n$?
\end{problem}

\begin{problem}
    \begin{enumerate}
        \item \textbf{\textcolor{red}{\heart}} Let $m$ and $n$ be nonnegative integers. Evaluate the integral
        \[
        B(m,n) = \int_0^1 x^m (1-x)^n dx,
        \]
        by interpreting the integral as a probability.
        \item Let $R$ be the region consisting of all triples $(x,y,z)$ of nonnegative real numbers satisfying $x+y+z\leq 1$. Let $w = 1-x-y-z$. Express the value of the triple integral (taken over the region $R$)
        \[
        \int\!\!\!\int\!\!\!\int x^1 y^9 z^8 w^4 dx\,dy\,dz
        \]
        in the form $a! b! c! d! / n!$, where $a$, $b$, $c$, $d$, and $n$ are positive integers.
    \end{enumerate}
\end{problem}

\begin{problem}
    \begin{enumerate}
        \item Choose $n$ points at random (uniformly and independently) on the circumference of a circle. Find the probability $p_n$ that all the points lie on a semicircle. (For instance, $p_1 = p_2 = 1$.)
        \item Fix $\theta < 2\pi$ and find the probability that the $n$ points lie on an arc subtending an angle $\theta$.
        \item Choose four points at random on the surface of a sphere. Find the probability that the center of the sphere is contained within the convex hull of the four points.
        \item (more difficult---I don’t know an elegant proof) Choose $n$ points uniformly at random in a square (or more generally, a parallelogram). Show that the probability that the points are in convex position (i.e., each is a vertex of their convex hull) is given by
        \[
        P_n = \frac{1}{n!} \binom{2n-2}{n-1}^2.
        \]
    \end{enumerate}
\end{problem}

\begin{problem}
Passengers $P_1,\ldots,P_n$ enter a plane with $n$ seats. Each passenger has a different assigned seat. The first passenger sits in the wrong seat. Thereafter, each passenger either sits in their seat if unoccupied or otherwise sits in a random unoccupied seat. What is the probability that the last passenger sits in his or her own seat?
\end{problem}

\begin{problem}
Let $x_1,x_2,\ldots,x_n$ be $n$ points (in that order) on the circumference of a circle. A person starts at the point $x_1$ and walks to one of the two neighboring points with probability $1/2$ for each. The person continues to walk in this way, always moving from the present point to one of the two neighboring points with probability $1/2$ for each. Find the probability $p_i$ that the point $x_i$ is the last of the $n$ points to be visited for the first time. In other words, find the probability that when $x_i$ is visited for the first time, all the other points will have already been visited. For instance, $p_1 = 0$ (when $n>1$), since $x_1$ is the first of the $n$ points to be visited.
\end{problem}

\begin{problem}
\textbf{\textcolor{red}{\heart}} There are $n$ parking spaces $1,2,\ldots,n$ (in that order) on a one-way street. Cars $C_1,\ldots,C_n$ enter the street in that order and try to park. Each car $C_i$ has a preferred space $a_i$. A car will drive to its preferred space and try to park there. If the space is already occupied, the car will park in the next available space. If the car must leave the street without parking, then the process fails. If $\alpha = (a_1,\ldots,a_n)$ is a sequence of preferences that allows every car to park, then we call $\alpha$ a parking function. For instance, there are $16$ parking functions of length $3$, given by (abbreviating $(1,1,1)$ as $111$, etc.) $111$, $112$, $121$, $211$, $113$, $131$, $311$, $122$, $212$, $221$, $123$, $132$, $213$, $231$, $312$, $321$. Show that the number of parking functions of length $n$ is equal to $(n+1)^{n-1}$.
\end{problem}

\begin{problem}
A snake on the $8 \times 8$ chessboard is a nonempty subset $S$ of the squares of the board obtained as follows: Start at one of the squares and continue walking one step up or to the right, stopping at any time. The squares visited are the squares of the snake. Here is an example of the $8 \times 8$ chessboard covered with disjoint snakes. Find the total number of ways to cover an $8 \times 8$ chessboard with disjoint snakes. Generalize to an $m \times n$ chessboard.
\end{problem}

\begin{problem}
Let $n$ balls be thrown uniformly and independently into $n$ bins.
    \begin{enumerate}
        \item Find the probability that bin 1 is empty.
        \item Find the expected number of empty bins.
    \end{enumerate}
\end{problem}

\begin{problem}
Consider $2m$ persons forming $m$ couples who live together at a given time. Suppose that at some later time, the probability of each person being alive is $p$, independently of other persons. At that later time, let $A$ be the number of persons that are alive and let $S$ be the number of couples in which both partners are alive. For any number of total surviving persons $a$, find expected number of surviving couples, i.e., $E[S~|~A = a]$.
\end{problem}

\begin{problem}
Let $X_1$, $X_2$, and $X_3$ be three independent, continuous random variables with the same distribution. Given that $X_2$ is smaller than $X_3$, what is the conditional probability that $X_1$ is smaller than $X_2$?
\end{problem}

\begin{problem}
    \begin{enumerate}
        \item In game show each contestant $i$ spins an infinitely calibrated ``fair'' wheel of fortune, which assigns the contestant a real number $X_i$ between 1 and 100. Find $P(X_1 < X_2)$.
        \item Find $P(X_1 < X_2, X_1 < X_3)$, i.e., find the probability that the first contestant will have the smallest value of the first three contestants.
        \item Consider a new random variable $N$, which is integer valued. $N$ is the index of the first contestant who is assigned a smaller number than contestant $1$. As an illustration, if contestant 1 has a smaller value than contestants $2$, $3$, and $4$, but contestant $5$ has a smaller value than contestant $1$ ($X_5 < X_1$), then $N = 5$. Find $P(N > n)$ as a function of $n$.
    \end{enumerate}
\end{problem}

\begin{problem}
The adventures of Ant Alice. On a meter-long rod sit $25$ ants placed uniformly at random, each facing a uniformly chosen direction (east or west). They proceed to march forward at $1$ cm/sec; whenever two ants collide, they reverse directions. Ants fall off the rod when they reach either endpoint. Alice is the ant initially positioned 13th from the west.
    \begin{enumerate}
        \item How long does it take before we can be certain that Alice is off the rod?
        \item What is the probability that Alice falls off the rod facing the same direction as her initial direction?
        \item What is the probability that Alice is the last ant to fall off the rod?
        \item What is the expected number of collisions that Alice has?
        \item What is the probability that Alice has more collisions than any other ant?
        \item Alice has a cold, which is transmitted from ant to ant instantly upon collision. How is the expected number of ants infected in the process?
    \end{enumerate}
\end{problem}



\section*{Inequalities}

\begin{problem}
For $p > 1$ and $a_1, a_2, \ldots, a_n$ positive, show that
\[
\sum_{k=1}^{n} \left( \frac{a_1 + a_2 + \cdots + a_k}{k} \right)^p
< \left( \frac{p}{p-1} \right)^p \sum_{k=1}^n a_k^p.
\]
\end{problem}

\begin{problem}
If $a_n > 0$ for $n = 1, 2, \ldots$, show that
\[
\sum_{n=1}^{\infty} \sqrt[n]{a_1 a_2 \cdots a_n} \leq e \sum_{n=1}^{\infty} a_n,
\]
provided that $\sum_{n=1}^{\infty} a_n$ converges.
\end{problem}

\begin{problem}
For $n = 1, 2, 3, \ldots$ let
\[
x_n = \frac{1000^n}{n!}.
\]
Find the largest term of the sequence.
\end{problem}

\begin{problem}
Suppose that $a_1, a_2, \ldots, a_n$ with $n \geq 2$ are real numbers greater than $-1$, and all the numbers $a_j$ have the same sign. Show that
\[
(1 + a_1)(1 + a_2) \cdots (1 + a_n) > 1 + a_1 + a_2 + \cdots + a_n.
\]
\end{problem}

\begin{problem}
If $a_1, \ldots, a_{n+1}$ are positive real numbers with $a_1 = a_{n+1}$, show that
\[
\prod_{i=1}^n \left(\frac{a_i}{a_{i+1}}\right)^n \geq \prod_{i=1}^n \frac{a_{i+1}}{a_i}.
\]
\end{problem}

\begin{problem}
Show that for any real numbers $a_1, a_2, \ldots, a_n$,
\[
\left(\sum_{i=1}^n \frac{a_i}{i}\right)^2
\leq
\sum_{i=1}^n \sum_{j=1}^n \frac{a_i a_j}{i + j - 1}.
\]
\end{problem}

\begin{problem}
Let $y = f(x)$ be a continuous, strictly increasing function of $x$ for $x \geq 0$, with $f(0) = 0$, and let $f^{-1}$ denote the inverse function to $f$. If $a$ and $b$ are nonnegative constants, then show that
\[
ab \leq \int_0^a f(x) dx + \int_0^b f^{-1}(y) dy.
\]
\end{problem}

\begin{problem}
\textbf{\textcolor{red}{\heart}}
Let $a_1, a_2, \ldots, a_n$ be real numbers. Show that
\[
\min_{i<j} (a_i - a_j)^2
\leq M_2 \left(
a_1^2 + \cdots + a_n^2
\right)
\]
where
\[
M_2 =
\frac{12}{n(n^2-1)}.
\]
\end{problem}

\begin{problem}
Let $f$ be a continuous function on the interval $[0, 1]$ such that $0 < m \leq f(x) \leq M$ for all $x$ in $[0, 1]$. Show that
\[
\left( \int_0^1 \frac{dx}{f(x)} \right) \left( \int_0^1 f(x) dx \right)
\leq \frac{(m + M)^2}{4mM}.
\]
\end{problem}

\begin{problem}
Consider any sequence $a_1, a_2, \ldots$ of real numbers. Show that
\[
\sum_{n=1}^{\infty} a_n \leq \frac{2}{\sqrt{3}} \sum_{n=1}^{\infty} \left( r_n n \right)^{1/2}
\]
where
\[
r_n = \sum_{k=n}^{\infty} a_k^2.
\]
(If the left-hand side of the inequality is $\infty$, then so is the right-hand side.)
\end{problem}

\begin{problem}
\textbf{\textcolor{red}{\heart}}
Show that
\[
\frac{1}{(n-1)!} \int_n^{\infty} w(t) e^{-t} dt < \frac{1}{(e-1)^n},
\]
where $t$ is real, $n$ is a positive integer, and
\[
w(t) = (t-1)(t-2) \cdots (t-n+1).
\]
\end{problem}

\begin{problem}
Let $(a_n)_{n=1}^{\infty}$ and $(b_n)_{n=1}^{\infty}$ be two sequences of positive numbers. Show that the following statements are equivalent:
\begin{itemize}
  \item There is a sequence $(c_n)_{n=1}^{\infty}$ of positive numbers such that $\sum_{n=1}^{\infty} \frac{a_n}{c_n}$ and $\sum_{n=1}^{\infty} \frac{c_n}{b_n}$ both converge.
  \item $\sum_{n=1}^{\infty} \sqrt{\frac{a_n}{b_n}}$ converges.
\end{itemize}
\end{problem}

\begin{problem}
Suppose that $a, b, c$ are real numbers in the interval $[-1, 1]$ such that $1 + 2abc \geq a^2 + b^2 + c^2$. Prove that $1 + 2(abc)^n \geq a^{2n} + b^{2n} + c^{2n}$ for all positive integers $n$.
\end{problem}

\begin{problem}
Suppose $f : \mathbb{R} \rightarrow \mathbb{R}$ is a two times differentiable function satisfying $f(0) = 1$, $f'(0) = 0$ and for all $x \in [0, \infty)$, it satisfies
\[
f''(x) - 5 f'(x) + 6 f(x) \geq 0
\]
Prove that, for all $x \in [0, \infty)$,
\[
f(x) \geq 3e^{2x} - 2e^{3x}
\]
\end{problem}

\begin{problem}
Let $f : [0, 1] \rightarrow \mathbb{R}$ be a continuous function satisfying $x f(y) + y f(x) \leq 1$ for every $x, y \in [0, 1]$.
\begin{enumerate}
    \item Show that
    \[
    \int_0^1 f(x) dx \leq \frac{\pi}{4}
    \]
    \item Find such a function for which equality occurs.
\end{enumerate}
\end{problem}

\begin{problem}
For what pairs of positive real numbers $(a, b)$ does the improper integral shown converge?
\[
\int_b^{\infty} \left( \sqrt{x+a} - \sqrt{x} - \sqrt{x} + \sqrt{x-b} \right) dx
\]
\end{problem}

\begin{problem}
Let $A$ be a positive real number. What are the possible values of
\[
\sum_{j=0}^{\infty} x_j^2,
\]
given that $x_0, x_1, \dots$ are positive numbers for which
\[
\sum_{j=0}^{\infty} x_j = A?
\]
\end{problem}

\begin{problem}
Let $f(x)$ be a continuous real-valued function defined on the interval $[0, 1]$. Show that
\[
\int_0^1 \int_0^1 |f(x) + f(y)| dx dy \geq \int_0^1 |f(x)| dx.
\]
\end{problem}

\begin{problem}
\textbf{\textcolor{red}{\heart}}
For each continuous function $f : [0, 1] \to \mathbb{R}$, let
\[
I(f) = \int_0^1 x^2 f(x) dx \quad \text{and} \quad J(f) = \int_0^1 x (f(x))^2 dx.
\]
Find the maximum value of $I(f) - J(f)$ over all such functions $f$.
\end{problem}

\begin{problem}
Suppose that $f : [0, 1] \to \mathbb{R}$ has a continuous derivative and that
\[
\int_0^1 f(x) dx = 0.
\]
Prove that for every $\alpha \in (0, 1)$,
\[
\left| \int_0^{\alpha} f(x) dx \right| \leq \frac{1}{8} \max_{0 \leq x \leq 1} |f'(x)|
\]
\end{problem}

\begin{problem}
For $m \geq 3$, a list of $\binom{m}{3}$ real numbers $a_{ijk}$ ($1 \leq i < j < k \leq m$) is said to be area definite for $\mathbb{R}^n$ if the inequality
\[
\sum_{1 \leq i < j < k \leq m} a_{ijk}\, \text{Area}(\triangle A_i A_j A_k) \geq 0
\]
holds for every choice of $m$ points $A_1, \ldots, A_m$ in $\mathbb{R}^n$. For example, the list of four numbers $a_{123} = a_{124} = a_{134} = 1$, $a_{234} = -1$ is area definite for $\mathbb{R}^2$. Prove that if a list of $\binom{m}{3}$ numbers is area definite for $\mathbb{R}^2$, then it is area definite for $\mathbb{R}^3$.
\end{problem}

\begin{problem}
Let $X_1, X_2, ...$ be independent random variables with the same distribution, and let $S_n = X_1 + X_2 + ... + X_n$, $n = 1, 2, ...$. For what real numbers $c$ is the following statement true:
\[
P\left( \left| \frac{S_{2n}}{2n} - c \right| \leq \left| \frac{S_n}{n} - c \right| \right) \geq \frac{1}{2}.
\]
\end{problem}

\begin{problem}
Let $H_k = \sum_{i=1}^{k} \frac{1}{i}$. Prove that the function $f : \mathbb{R} \to \mathbb{R}$ defined by
\[
f(x) = 1 + \sum_{n=1}^{\infty} \frac{x^n}{\prod_{k=1}^{n} H_k}
\]
has no real zeros.
\end{problem}

\begin{problem}
Let $f$ be a continuous, nonnegative function on $[0, 1]$. Show that
\[
\int_0^1 f(x)^3 dx \geq 4 \left( \int_0^1 x f(x)^2 dx \right) \left( \int_0^1 x^2 f(x) dx \right)
\]
\end{problem}

\begin{problem}
Let $f : \mathbb{R} \to \mathbb{R}$ be an infinitely differentiable function satisfying $f(0) = 0$, $f(1) = 1$, and $f(x) \geq 0$ for all $x \in \mathbb{R}$. Show that there exists a positive integer $n$ and a real number $x$ such that $f^{(n)}(x) < 0$.
\end{problem}

\begin{problem}
Let $f = (f_1, f_2)$ be a function from $\mathbb{R}^2$ to $\mathbb{R}^2$ with continuous partial derivatives $\frac{\partial f_i}{\partial x_j}$ that are positive everywhere. Suppose that
\[
\frac{\partial f_1}{\partial x_1} \frac{\partial f_2}{\partial x_2}
- \frac{1}{4}
\left(\frac{\partial f_1}{\partial x_2} + \frac{\partial f_2}{\partial x_1}\right)^2
> 0
\]
everywhere. Prove that $f$ is one-to-one.
\end{problem}

\begin{problem}
Determine the greatest possible value of $\sum_{i=1}^{10} \cos(3x_i)$ for real numbers $x_1, x_2, \ldots, x_{10}$ satisfying $\sum_{i=1}^{10} \cos(x_i) = 0$.
\end{problem}

\begin{problem}
Let $S_a = \{ f : [-a, a] \to \mathbb{R} \mid \int_{-a}^a f(x)^2 dx = 1 \}$ and define $M_j(f) = \int_{-a}^a x^j f(x)^2 dx$.
\begin{itemize}
    \item What is $\sup (M_0(f)^2 + M_1(f)^2)$?
    \item What is $\sup (M_0(f)^2 + M_1(f)^2 + M_2(f)^2)$?
\end{itemize}
\end{problem}

\begin{problem}
Let $n$ be a sufficiently large prime and $f : \mathbb{Z}/n\mathbb{Z} \to \mathbb{C}$ be a function such that $|f(x)| \leq 1$ for all $x \in \mathbb{Z}/n\mathbb{Z}$. Prove that
\[
\mathbb{E}_{x, y \in \mathbb{Z}/n\mathbb{Z}}[f(x) f(x+y) f(x+2y) f(x+3y)]
\leq
\left(
\mathbb{E}_{x, h_1, h_2, h_3 \in \mathbb{Z}/n\mathbb{Z}}[
    f(x) \overline{f(x + h_1)} \overline{f(x + h_2)} \overline{f(x + h_3)}
    f(x + h_1 + h_2) f(x + h_1 + h_3) f(x + h_2 + h_3)
    \overline{f(x + h_1 + h_2 + h_3)}
]
\right)^{1/8}
\]
When does equality occur?
\end{problem}

\begin{problem}
Given a random variable $X$ taking values in $\mathbb{R}$, define $g(\Theta) = \mathbb{E}[e^{i \Theta X}]$. Prove that
\[
P[|X| \leq 1] \leq 100000 \int_{-1}^1 |g(\Theta)| d\Theta.
\]
\end{problem}


\section*{Combinatorial Configurations}

These problems aim to capture the flavors of combinatorial questions which have been common in recent years on the Putnam. Difficulty is very roughly increasing, with the first several being more exercise-like.

\begin{problem}
Can a $2018 \times 2018$ grid be tiled with rotations and reflections of the L-tetromino?
\end{problem}

\begin{problem}
Each edge of a complete graph on $2k$ vertices is removed with probability $\frac{1}{2}$. Prove that with probability greater than $\frac{1}{4k}$, the maximum degree of the remaining graph is at most $k-1$.
\end{problem}

\begin{problem}
Let $n$ be a fixed positive integer. How many ways are there to write $n$ as a sum of positive integers,
\[
n = a_1 + a_2 + \cdots + a_k,
\]
with $k$ an arbitrary positive integer and $a_1 \leq a_2 \leq \cdots \leq a_k \leq a_1 + 1$?
\end{problem}

\begin{problem}
How many $n$-digit numbers whose digits are in the set $\{2, 3, 7, 9\}$ are divisible by 3?
\end{problem}

\begin{problem}
A $6 \times 6$ grid is tiled by dominoes. Prove that there exists some line which cuts the board into two nonempty parts without cutting through any domino.
\end{problem}

\begin{problem}
\textbf{\textcolor{red}{\heart}}
A collection of $2020$ sets, each of size $42$, has the property that any two of the sets have exactly one common element. Must all of the sets have a common element?
\end{problem}

\begin{problem}
On each face of a regular icosahedron is written a nonnegative integer such that the sum of all $20$ integers is $39$. Show that there are two faces that share a vertex and have the same integer written on them.
\end{problem}

\begin{problem}
Given a positive integer $n$, what is the largest $k$ such that the numbers $1, 2, \ldots, n$ can be put into $k$ boxes such that the sum of the numbers in each box is the same?
\end{problem}

\begin{problem}
\textbf{\textcolor{red}{\heart}}
A class with $2N$ students took a quiz, where each student received an integer score between $0$ and $10$ inclusive. At least one student received each possible score, and the average score was exactly $7.4$. Show that the class can be divided into two groups with equal number of students and the same average score.
\end{problem}

\begin{problem}
Callie and Marie play a game in which they take turns filling entries of an initially empty $2008 \times 2008$ array. Callie plays first. At each turn, a player chooses a real number and places it in a vacant entry. The game ends when all the entries are filled. Callie wins if the determinant of the resulting matrix is nonzero; Marie wins if it is zero. Which player has a winning strategy?
\end{problem}

\begin{problem}
Consider a $(2m-1) \times (2n-1)$ rectangular region, where $m, n$ are integers such that $m, n \geq 4$. This region is to be tiled without overlap using rotations and reflections of the bent triomino and the S-tetromino. What is the fewest number of these tiles which can be used?
\end{problem}

\begin{problem}
\textbf{\textcolor{red}{\heart}}
Let $B$ be a set of more than $2^{n+1}/n$ distinct points with coordinates of the form $\{\pm 1, \ldots, \pm 1\}$ in $n$-dimensional space with $n \geq 3$. Show that there are three distinct points in $B$ which are the vertices of an equilateral triangle.
\end{problem}

\begin{problem}
A set of $n$ points is given in the plane such that the distance between any two of them is greater than $1$. Prove that one can choose at least $\frac{n}{7}$ of these points such that the distance between any two of them is greater than $\sqrt{3}$.
\end{problem}

\begin{problem}
A round-robin tournament of $2n$ teams lasted for $2n-1$ days, as follows. On each day, every team played one game against another team, with one team winning and one team losing in each of the $n$ games. Over the course of the tournament, each team played every other team exactly once. Can one necessarily choose one winning team from each day without choosing any team more than once?
\end{problem}

\begin{problem}
Callie and Marie play a game with $r$ red and $g$ green stones. At every step, one can remove an amount $k$ of stones of either color, where $k$ divides the number of stones remaining of the other color. She who moves last wins, and Callie starts. For which $(r, g)$ does Callie win?
\end{problem}

\begin{problem}
Consider a polyhedron with at least five faces such that exactly three edges emerge from each of its vertices, and no two edges share more than one face. Two players play the following game:\\
Each player, in turn, signs his or her name on a previously unsigned face. The winner is the player who first succeeds in signing three faces that share a common vertex.\\
Show that the player who signs first will always win by playing as well as possible.
\end{problem}

\begin{problem}
Two hundred students participated in a mathematical contest. They had six problems to solve. It is known that each problem was correctly solved by at least $120$ participants. Prove that there must be two participants such that every problem was solved by at least one of these two students.
\end{problem}

\begin{problem}
A city is a point on the plane. Suppose there are $n \geq 2$ cities. Suppose that for each city $X$, there is another city $N(X)$ that is strictly closer to $X$ than all the other cities. The government builds a road connecting each city $X$ and its $N(X)$; no other roads have been built. Suppose we know that, starting from any city, we can reach any other city through a series of roads.\\
We call a city $Y$ \emph{suburban} if it is $N(X)$ for some city $X$. Show that there are at least $\frac{n-2}{4}$ suburban cities.
\end{problem}

\begin{problem}
Let $n \geq 2$ be an integer and $T_n$ be the number of nonempty subsets $S$ of $\{1, 2, 3, \ldots, n\}$ with the property that the average of the elements in $S$ is an integer. Prove that $T_n - n$ is always even.
\end{problem}

\begin{problem}
Let $\mathbb{Z}^n$ be the integer lattice in $\mathbb{R}^n$. Two points in $\mathbb{Z}^n$ are called neighbors if they differ by exactly $1$ in one coordinate and are equal in all other coordinates. For which integers $n \geq 1$ does there exist a set of points $S \subset \mathbb{Z}^n$ satisfying the following two conditions?
\begin{enumerate}
    \item If $p$ is in $S$, then none of the neighbors of $p$ is in $S$.
    \item If $p \in \mathbb{Z}^n$ is not in $S$, then exactly one of the neighbors of $p$ is in $S$.
\end{enumerate}
\end{problem}

\begin{problem}
Prove that it is not possible to color the squares of a $11 \times 11$ grid using three colors such that no four squares whose centers form the vertices of a rectangle with sides parallel to the sides of the grid, have the same color.
\end{problem}

\begin{problem}
A circle is divided by $2n$ evenly spaced points into $2n$ equal arcs. Half of these arcs are colored blue and half are colored red. The blue arcs are numbered from $1$ to $n$ in the counterclockwise direction, starting from an arbitrary blue arc, and red arcs are numbered from $1$ to $n$ in the clockwise direction, starting from an arbitrary red arc. Prove that one can find $n$ consecutive arcs, each of which has a different label.
\end{problem}

\begin{problem}
Let $n$ and $k$ be positive integers. Cathy is playing the following game. There are $n$ marbles and $k$ boxes, with the marbles labelled $1$ to $n$. Initially, all marbles are placed inside one box. Each turn, Cathy chooses a box and then moves the marbles with the smallest label, say $i$, to either any empty box or the box containing marble $i+1$. Cathy wins if at any point there is a box containing only marble $n$. Determine all pairs of integers $(n, k)$ such that Cathy can win this game.
\end{problem}

\begin{problem}
Callie and Marie play a game in which they take turns choosing an element from the group of invertible $n \times n$ matrices with entries in the field $\mathbb{Z}/p\mathbb{Z}$ of integers modulo $p$, where $n$ is fixed positive integer and $p$ is a fixed prime number. The rules of the game are:
\begin{itemize}
    \item A player cannot choose an element that has been chosen by either player on any previous turn.
    \item A player can only choose an element that commutes with all previously chosen elements.
    \item A player who cannot choose an element on his/her turn loses the game.
\end{itemize}
Callie takes the first turn. For which $(n, p)$ does Callie have a winning strategy?
\end{problem}

\begin{problem}
Suppose a finite number of integer arithmetic progressions are given, such that every positive integer belongs to exactly one of them. Prove that two of these progressions have the same difference.
\end{problem}

\begin{problem}
Suppose you are given a binary string which is not entirely zeroes. Prove that you may insert $+$ between some of the digits so that the resulting summation, when carried out in base 2, yields a power of $2$. Example: for $101111$ the split $1 + 0 + 1111$ is possible.
\end{problem}

\begin{problem}
An alphabet consists of three letters. Some of the words of length $2$ or more are prohibited, and all of the prohibited sequences have different lengths. A word is a sequence of letters of any length. A correct word does not contain any prohibited sequence. Prove that there are correct words of any length.
\end{problem}

\begin{problem}
\textbf{\textcolor{red}{\heart}}
There are $2007$ senators in a senate. Certain pairs of senators in the senate are enemies with each other. Prove that there is a non-empty subset $K$ of senators such that, for every senator in the senate, the number of enemies of that senator in the set $K$ is an even number.
\end{problem}

\begin{problem}
Call a subset $S$ of $\{1, 2, \ldots, n\}$ mediocre if it has the following property: Whenever $a, b$ are elements of $S$ whose average is an integer, that average is also an element of $S$. Let $A(n)$ be the number of mediocre subsets of $\{1, 2, \ldots, n\}$. Find all positive integers $n$ such that $A(n + 2) + A(n) = 2A(n + 1) + 1$.
\end{problem}

\begin{problem}
Let $S_1, S_2, \ldots, S_m$ be distinct subsets of $\{1, 2, \ldots, n\}$ such that $|S_i \cap S_j| = 1$ for all $i \neq j$. Prove that $m \leq n$.
\end{problem}

\begin{problem}
\textbf{\textcolor{red}{\heart}}
For a set $S$ of nonnegative integers, let $r_S(n)$ denote the number of ordered pairs $(s_1, s_2)$ such that $s_1 \in S$, $s_2 \in S$, $s_1 \neq s_2$, $s_1 + s_2 = n$. Is it possible to partition the nonnegative integers into two sets $A, B$ in such a way that $r_A(n) = r_B(n)$ for all $n$?
\end{problem}

\begin{problem}
The $30$ edges of a regular icosahedron are distinguished by labeling them $1, 2, \ldots, 30$. How many different ways are there to paint each edge red, white, or blue such that each of the $20$ triangular faces of the icosahedron has two edges of the same color and a third edge of a different color?
\end{problem}

\begin{problem}
There are $2010$ boxes labeled $B_1, B_2, \ldots, B_{2010}$ and $2010n$ balls have been distributed among them, for some positive integer $n$. You may redistribute the balls by a sequence of moves, each of which consists of choosing an $i$ and moving exactly $i$ balls from box $B_i$ into any one other box. For which values of $n$ is it possible to reach the distribution with exactly $n$ balls in each box, regardless of the initial distribution of balls?
\end{problem}

\begin{problem}
Let $f(n)$ be the minimum number of cells required to be selected in an $n \times n$ grid so that any L-tetromino contains at least one selected cell. Prove that there exists some constant $c$ for which $cn^2 - 10n \leq f(n) \leq cn^2 + 10n$ for all $n$.\\
(Bonus: Can you prove this for any shape?)
\end{problem}

\section*{Congruences and Divisibility}

\begin{problem}
Let $n_1, n_2, \ldots, n_s$ be distinct integers such that
\[
(n_1 + k)(n_2 + k) \cdots (n_s + k)
\]
is an integral multiple of $n_1 n_2 \cdots n_s$ for every integer $k$. For each of the following assertions, give a proof or a counterexample:
\begin{itemize}
    \item $|n_i| = 1$ for some $i$.
    \item If further all $n_i$ are positive, then $\{n_1, n_2, \ldots, n_s\} = \{1, 2, \ldots, s\}$.
\end{itemize}
\end{problem}

\begin{problem}
How many coefficients of the polynomial
\[
P_n(x_1, \ldots, x_n) = \prod_{1 \leq i < j \leq n} (x_i + x_j)
\]
are odd?
\end{problem}

\begin{problem}
If $p$ is a prime number greater than $3$ and $k = \lfloor 2p/3 \rfloor$, prove that the sum
\[
\binom{p}{1} + \binom{p}{2} + \cdots + \binom{p}{k}
\]
of binomial coefficients is divisible by $p^2$.
\end{problem}

\begin{problem}
Do there exist positive integers $a$ and $b$ with $b - a > 1$ such that for every $a < k < b$, either $\gcd(a, k) > 1$ or $\gcd(b, k) > 1$?
\end{problem}

\begin{problem}
Suppose that $f(x)$ and $g(x)$ are polynomials (with $f(x)$ not identically $0$) taking integers to integers such that for all $n \in \mathbb{Z}$, either $f(n) = 0$ or $f(n) \mid g(n)$. Show that $f(x) \mid g(x)$, i.e., there is a polynomial $h(x)$ with rational coefficients such that $g(x) = f(x) h(x)$.
\end{problem}

\begin{problem}
Let $q$ be an odd positive integer, and let $N_q$ denote the number of integers $a$ such that $0 < a < q/4$ and $\gcd(a, q) = 1$. Show that $N_q$ is odd if and only if $q$ is of the form $p^k$ with $k$ a positive integer and $p$ a prime congruent to 5 or 7 modulo 8.
\end{problem}

\begin{problem}
Let $p$ be in the set $\{3, 5, 7, 11, \ldots\}$ of odd primes, and let
\[
F(n) = 1 + 2n + 3n^2 + \cdots + (p-1) n^{p-2}.
\]
Prove that if $a$ and $b$ are distinct integers in $\{0, 1, 2, \ldots, p-1\}$ then $F(a)$ and $F(b)$ are not congruent modulo $p$, that is, $F(a) - F(b)$ is not exactly divisible by $p$.
\end{problem}

\begin{problem}
Do there exist $1,000,000$ consecutive integers each of which contains a repeated prime factor?
\end{problem}

\begin{problem}
\textbf{\textcolor{red}{\heart}}
A positive integer $n$ is powerful if for every prime $p$ dividing $n$, we have that $p^2$ divides $n$. Show that for any $k \geq 1$ there exist $k$ consecutive integers, none of which is powerful.
\end{problem}

\begin{problem}
Show that for any $k \geq 1$ there exist $k$ consecutive positive integers, none of which is a sum of two squares. (You may use the fact that a positive integer $n$ is a sum of two squares if and only if for every prime $p \equiv 3 \pmod 4$, the largest power of $p$ dividing $n$ is an even power of $p$.)
\end{problem}

\begin{problem}
Prove that every positive integer has a multiple whose decimal representation involves all ten digits.
\end{problem}

\begin{problem}
Prove that among any ten consecutive integers at least one is relatively prime to each of the others.
\end{problem}

\begin{problem}
Find the length of the longest sequence of equal nonzero digits in which an integral square can terminate (in base $10$), and find the smallest square which terminates in such a sequence.
\end{problem}

\begin{problem}
Show that if $n$ is an odd integer greater than $1$, then $n$ does not divide $2^n + 2$.
\end{problem}

\begin{problem}
For positive integer $a$, we define the series
\[
f_a(q) = \sum_{k \geq 0,\, a k + 1\ \text{is a square}} q^k.
\]
Find all positive integer triples $(a, b, c)$ such that
\[
f_a(q) \equiv f_b(q) f_c(q) \pmod{2}
\]
which means that the corresponding coefficients match modulo $2$. (Hint: Use a computer to find a few triples, then look for patterns.)
\end{problem}

\begin{problem}
Define a sequence $\{a_i\}$ by $a_1 = 3$ and $a_{i+1} = 3 a_i$ for $i \geq 1$. Which integers between $00$ and $99$ inclusive occur as the last two digits in the decimal expansion of infinitely many $a_i$?
\end{problem}

\begin{problem}
What is the units (i.e., rightmost) digit of
\[
\left\lfloor \frac{10^{20000}}{10^{100} + 3} \right\rfloor?
\]
Here $[x]$ is the greatest integer $\leq x$.
\end{problem}

\begin{problem}
Suppose $p$ is an odd prime. Prove that
\[
\sum_{j = 0}^p \binom{p}{j} \binom{p+j}{j} \equiv 2p + 1 \pmod{p^2}.
\]
\end{problem}

\begin{problem}
Prove that for $n \geq 2$,
\[
\underbrace{2^{2 \cdots 2}}_{n\ \text{terms}} \equiv \underbrace{2^{2 \cdots 2}}_{n-1\ \text{terms}} \pmod{n}.
\]
\end{problem}

\begin{problem}
The sequence $(a_n)_{n \geq 1}$ is defined by $a_1 = 1$, $a_2 = 2$, $a_3 = 24$, and, for $n \geq 4$,
\[
a_n = \frac{6a_{n-1}^2 a_{n-3} - 8 a_{n-1} a_{n-2}^2}{a_{n-2} a_{n-3}}.
\]
Show that, for all $n$, $a_n$ is an integer multiple of $n$.
\end{problem}

\begin{problem}
\textbf{\textcolor{red}{\heart}}
Prove that the expression
\[
\frac{\gcd(m, n)}{n} \binom{n}{m}
\]
is an integer for all pairs of integers $n \geq m \geq 1$.
\end{problem}

\begin{problem}
Show that for each positive integer $n$,
\[
n! = \prod_{i = 1}^n \mathrm{lcm}\{1, 2, \ldots, \lfloor n/i \rfloor\}.
\]
(Here lcm denotes the least common multiple, and $\lfloor x \rfloor$ denotes the greatest integer $\leq x$.)
\end{problem}

\begin{problem}
Define a sequence $\{u_n\}_{n=0}^{\infty}$ by $u_0 = u_1 = u_2 = 1$, and thereafter by the condition that
\[
\det\begin{pmatrix}
u_n & u_{n+1} \\
u_{n+2} & u_{n+3}
\end{pmatrix} = n!
\]
for all $n \geq 0$. Show that $u_n$ is an integer for all $n$. (By convention, $0! = 1$.)
\end{problem}

\begin{problem}
\textbf{\textcolor{red}{\heart}}
Let $p$ be a prime number. Let $h(x)$ be a polynomial with integer coefficients such that $h(0), h(1), \ldots, h(p^2-1)$ are distinct modulo $p^2$. Show that $h(0), h(1), \ldots, h(p^3-1)$ are distinct modulo $p^3$.
\end{problem}

\begin{problem}
Let $f(x) = a_0 + a_1 x + \cdots$ be a power series with integer coefficients, with $a_0 \neq 0$. Suppose that the power series expansion of $\frac{f'(x)}{f(x)}$ at $x = 0$ also has integer coefficients. Prove or disprove that $a_0 \mid a_n$ for all $n \geq 0$.
\end{problem}

\begin{problem}
Let $S$ be a set of rational numbers such that
\begin{enumerate}
    \item $0 \in S$,
    \item If $x \in S$ then $x + 1 \in S$ and $x - 1 \in S$, and
    \item If $x \in S$ and $x \notin \{0, 1\}$, then $1/(x(x - 1)) \in S$.
\end{enumerate}
Must $S$ contain all rational numbers?
\end{problem}

\begin{problem}
\textbf{\textcolor{red}{\heart}}
Prove that for each positive integer $n$, the number $10^{10^{n}} + 10^{10^n} + 10^n - 1$ is not prime.
\end{problem}

\begin{problem}
Let $p$ be an odd prime. Show that for at least $(p+1)/2$ values of $n$ in $\{0, 1, 2, \ldots, p-1\}$,
\[
\sum_{k=0}^{p-1} k! n^k
\]
is not divisible by $p$.
\end{problem}

\begin{problem}
Let $a$ and $b$ be distinct rational numbers such that $a^n - b^n$ is an integer for all positive integers $n$. Prove or disprove that $a$ and $b$ must themselves be integers.
\end{problem}

\begin{problem}
Find the smallest integer $n \geq 2$ for which there exists an integer $m$ with the following property: for each $i \in \{1, \ldots, n\}$, there exists $j \in \{1, \ldots, n\}$ different from $i$ such that $\gcd(m + i, m + j) > 1$.
\end{problem}

\begin{problem}
Let $p$ be an odd prime number such that $p \equiv 2 \pmod 3$. Define a permutation $\pi$ of the residue classes modulo $p$ by $\pi(x) \equiv x^3 \pmod p$. Show that $\pi$ is an even permutation if and only if $p \equiv 3 \pmod 4$.
\end{problem}

\begin{problem}
Suppose that a positive integer $N$ can be expressed as the sum of $k$ consecutive positive integers
\[
N = a + (a+1) + (a+2) + \cdots + (a + k - 1)
\]
for $k = 2017$ but for no other values of $k > 1$. Considering all positive integers $N$ with this property, what is the smallest positive integer $a$ that occurs in any of these expressions?
\end{problem}

\begin{problem}
Let $n$ be a positive integer. Prove that
\[
\sum_{k = 1}^n (-1)^{\lfloor k(\sqrt{2} - 1) \rfloor} \geq 0.
\]
\end{problem}

\begin{problem}
How many positive integers $N$ satisfy all of the following three conditions? (i) $N$ is divisible by $2020$; (ii) $N$ has at most $2020$ decimal digits; (iii) The decimal digits of $N$ are a string of consecutive ones followed by a string of consecutive zeros.
\end{problem}

\begin{problem}
Find all positive integers $n < 10^{100}$ for which simultaneously $n \mid 2^n$, $n-1 \mid 2^{n-1}$, and $n-2 \mid 2^{n-2}$.
\end{problem}

\begin{problem}
Define the Fibonacci numbers by $F_1 = F_2 = 1$ and $F_n = F_{n-1} + F_{n-2}$ for $n \geq 3$. Let $k$ be a positive integer. Suppose that for every positive integer $m$ there exists a positive integer $n$ such that $m \mid F_n - k$. Must $k$ be a Fibonacci number?
\end{problem}

\begin{problem}
Suppose that
\[
f(x) = \sum_{i = 0}^{\infty} c_i x^i
\]
is a power series for which each coefficient $c_i$ is $0$ or $1$. Show that if $f(2/3) = 3/2$, then $f(1/2)$ must be irrational.
\end{problem}

\begin{problem}
Let $P(x)$ be a polynomial with integer coefficients such that $n$ divides $P(2n)$ for all positive integers $n$. Prove that $P$ is the zero polynomial.
\end{problem}

\begin{problem}
\textbf{\textcolor{red}{\heart}}
Let $a_1, a_2, \ldots, a_n$ be positive integers with product $P$, where $n$ is an odd positive integer. Prove that
\[
\gcd(a_1^n + P, a_2^n + P, \ldots, a_n^n + P) \leq 2\, \gcd(a_1, \ldots, a_n)^n.
\]
\end{problem}

\begin{problem}
Let $N$ denote the set of positive integers. Find all functions $f : N \to N$ such that
\[
n + f(m) \mid f(n) + n f(m)
\]
for all $m, n \in N$.
\end{problem}

\begin{problem}
Determine all positive integers $M$ such that the sequence $a_0, a_1, a_2, \ldots$ defined by
\[
a_0 = \frac{M + 1}{2}, \quad a_{k+1} = a_k \lfloor a_k \rfloor
\]
for $k = 0, 1, 2, \ldots$ contains at least one integer term.
\end{problem}

\begin{problem}
Find all positive integers $n > 2$ such that
\[
n!\ \Bigg|\ \sum_{\substack{p < q \leq n \\ p, q\ \text{primes}}} (p + q).
\]
\end{problem}

\begin{problem}
Let $a_0 = 1$, $a_1 = 2$, and $a_n = 4 a_{n-1} - a_{n-2}$ for $n \geq 2$. Find an odd prime factor of $a_{2015}$.
\end{problem}


\section*{Generating Functions}

\begin{problem}
Let $a_n = (n^2 + 1)3^n$. Compute $y = \sum_{n \geq 0} a_n x^n$.
\end{problem}

\begin{problem}
Compute
\[
y = x + \frac{2}{3}x^3 + \frac{2}{3}\frac{4}{5}x^5 + \frac{2}{3}\frac{4}{5}\frac{6}{7}x^7 + \cdots.
\]
\end{problem}

\begin{problem}
Given $a_0 = 2$, $a_1 = 3$, and
\[
(n+1)(n+2)a_{n+2} - 3(n+1)a_{n+1} + 2a_n = 0, \quad n \geq 0,
\]
compute $y = \sum_{n \geq 0} a_n x^n$.
\end{problem}

\begin{problem}
Given $a_0 = 1$ and $a_{n+1} = (n+1)a_n - \binom{n}{2} a_{n-2}$ for $n \geq 0$, compute
\[
y = \sum_{n \geq 0} \frac{a_n}{n!} x^n.
\]
\end{problem}

\begin{problem}
Let $k$ be a positive integer and let $m = 6k - 1$. Define
\[
S(m) = \sum_{j=1}^{2k-1} (-1)^{j+1} \binom{m}{3j-1}.
\]
For example with $k=3$,
\[
S(17) = \binom{17}{2} - \binom{17}{5} + \binom{17}{8} - \binom{17}{11} + \binom{17}{14}.
\]
Prove that $S(m)$ is never zero.
\end{problem}

\begin{problem}
For nonnegative integers $n$ and $k$, define $Q(n,k)$ to be the coefficient of $x^k$ in the expansion of $(1+x+x^2+x^3)^n$. Prove that
\[
Q(n,k) = \sum_{j=0}^n \binom{n}{j}\binom{n}{k-2j}.
\]
\end{problem}

\begin{problem}
Let $a_{m,n}$ denote the coefficient of $x^n$ in the expansion of $(1+x+x^2)^m$. Prove that for all $k \geq 0$,
\[
0 \leq \sum_{i=0}^{\lfloor 2k/3 \rfloor} (-1)^i a_{k-i,i} \leq 1.
\]
\end{problem}

\begin{problem}
Consider the power series expansion
\[
\frac{1}{1-2x-x^2} = \sum_{n=0}^\infty a_n x^n.
\]
Prove that, for each $n \geq 0$, there exists $m$ such that $a_{2n} + a_{2n+1} = a_m$.
\end{problem}

\begin{problem}
Let $A = \{(x,y): 0 \leq x,y \leq 1\}$. For $(x,y) \in A$, let
\[
S(x,y) = \sum_{\substack{1/2 \leq m/n \leq 2 \\ m,n \in \mathbb{Z}^+}} x^m y^n.
\]
Evaluate
\[
\lim_{\substack{(x,y) \to (1,1) \\ (x,y) \in A}} (1-xy^2)(1-x^2y) S(x,y).
\]
\end{problem}

\begin{problem}
For a set $S$ of nonnegative integers, let $r_S(n)$ denote the number of ordered pairs $(s_1,s_2)$ such that $s_1,s_2 \in S$, $s_1 \neq s_2$, and $s_1+s_2=n$. Is it possible to partition $\mathbb{N}$ into two sets $A,B$ such that $r_A(n)=r_B(n)$ for all $n$?
\end{problem}

\begin{problem}
For positive integers $m,n$, let $f(m,n)$ denote the number of $n$-tuples $(x_1,\dots,x_n)$ of integers such that $|x_1|+\cdots+|x_n| \leq m$. Show that $f(m,n)=f(n,m)$.
\end{problem}

\begin{problem}
Let $S_n$ be the set of permutations of $\{1,2,\dots,n\}$. For $\pi \in S_n$, let $\sigma(\pi)=1$ if $\pi$ is even, $-1$ if odd, and let $\nu(\pi)$ be the number of fixed points. Find
\[
\sum_{\pi \in S_n} \frac{\sigma(\pi)}{\nu(\pi)+1}.
\]
\end{problem}

\begin{problem}
Let $S$ be the set of sequences of length $2018$ with terms in $\{1,2,3,4,5,6,10\}$ summing to $3860$. Prove that
\[
|S| \leq 2^{3860} \binom{2018}{2048}^{2018}.
\]
\end{problem}

\begin{problem}
Suppose $\mathbb{Z}$ is written as a disjoint union of $n < \infty$ arithmetic progressions with differences $d_1 \geq d_2 \geq \cdots \geq d_n \geq 1$. Show that $d_1 = d_2$.
\end{problem}

\begin{problem}
Solve for $F(x,y) = \sum_{m,n \geq 0} a_{mn} x^m y^n$, where $a_{mn} \in \mathbb{R}$:
\[
(xy^2 + x - y)F(x,y) = xF(x,0) - y.
\]
The solution should admit a power series expansion at $(0,0)$.
\end{problem}

\begin{problem}
Find a simple description of coefficients $a_n \in \mathbb{Z}$ of the power series
\[
F(x) = x + a_2 x^2 + a_3 x^3 + \cdots
\]
satisfying
\[
F(x) = (1+x)F(x^2) + \frac{x}{1-x^2}.
\]
\end{problem}

\begin{problem}
Consider
\[
y = \sum_{n=0}^\infty \binom{3n}{n} x^n = 1 + 3x + 15x^2 + \cdots.
\]
Show that
\[
(27x-4)y^3 + 3y + 1 = 0.
\]
\end{problem}

\begin{problem}
Find the unique power series
\[
y = 1 + \tfrac{1}{2}x + \tfrac{1}{12}x^2 - \tfrac{1}{720}x^4 + \tfrac{1}{30240}x^6 + \cdots
\]
such that for all $n \geq 0$, the coefficient of $x^n$ in $y^{n+1}$ is $1$.
\end{problem}

\begin{problem}
Find the unique power series
\[
y = 1 + x - \tfrac{1}{2}x^2 + \tfrac{2}{3}x^3 + \cdots
\]
such that the constant term is $1$, the coefficient of $x$ is $1$, and for all $n \geq 2$ the coefficient of $x^n$ in $y^n$ is $0$.
\end{problem}

\begin{problem}
Let $f(m,0)=f(0,n)=1$ and for $m,n > 0$ let
\[
f(m,n) = f(m-1,n) + f(m,n-1) + f(m-1,n-1).
\]
Show that
\[
\sum_{n=0}^\infty f(n,n)x^n = \frac{1}{\sqrt{1-6x+x^2}}.
\]
\end{problem}

\section*{Polynomials}

\begin{problem}
Find the cubic equation whose roots are the cubes of the roots of
\[
x^3 + ax^2 + bx + c = 0.
\]
\end{problem}

\begin{problem}
(a) Determine all rational values for which $a, b, c$ are the roots of
\[
x^3 + ax^2 + bx + c = 0.
\]
(b) Show that the only real polynomials $\prod_{i=0}^{n-1} (x-a_i) = x^n + a_{n-1}x^{n-1} + \cdots + a_0$ in addition to those given by (a) are $x^n$, $x^2+x-2$, and exactly two others, approximately
\begin{align*}
x^3 + 0.56519772x^2 - 1.76929234x + 0.63889690, \\
x^4 + x^3 - 1.7548782x^2 - 0.5698401x + 0.3247183.
\end{align*}
\end{problem}

\begin{problem}
Assuming that all the roots of $x^3 + ax^2 + bx + c$ are real, show that the difference between the greatest and least roots is not less than $\sqrt{a^2-3b}$ nor greater than $2\sqrt{(a^2-3b)/3}$.
\end{problem}

\begin{problem}
The nonconstant polynomials $P(z)$ and $Q(z)$ with complex coefficients have the same set of zeros but possibly different multiplicities. The same is true for $P(z)+1$ and $Q(z)+1$. Prove that $P(z)=Q(z)$.
\end{problem}

\begin{problem}
If $a_0, a_1, \dots, a_n$ are real numbers satisfying
\[
\frac{a_0}{1} + \frac{a_1}{2} + \cdots + \frac{a_n}{n+1} = 0,
\]
show that $a_0 + a_1x + \cdots + a_nx^n = 0$ has at least one real root.
\end{problem}

\begin{problem}
Determine all polynomials of the form $\sum_{i=0}^n a_ix^{n-i}$ with $a_i = \pm 1$ such that each has only real zeros.
\end{problem}

\begin{problem}
Let $P(x)$ be a polynomial with real coefficients and form
\[
Q(x) = (x^2+1)P(x)P'(x) + x(P(x)^2 + P'(x)^2).
\]
If $P(x)=0$ has $n$ distinct real roots exceeding $1$, prove or disprove that $Q(x)=0$ has at least $2n-1$ distinct real roots.
\end{problem}

\begin{problem}
Prove that if
\[
11z^{10} + 10iz^9 + 10iz - 11 = 0,
\]
then $|z|=1$.
\end{problem}

\begin{problem}
Is there an infinite sequence $a_0,a_1,a_2,\dots$ of nonzero reals such that for each $n\geq 1$, the polynomial
\[
p_n(x) = a_0 + a_1x + \cdots + a_nx^n
\]
has exactly $n$ distinct real roots?
\end{problem}

\begin{problem}
Find all real polynomials $p(x)$ of degree $n\geq 2$ such that there exist $r_1<r_2<\cdots<r_n$ with
\begin{align*}
p(r_i)&=0, \quad i=1,\dots,n,\\
p'\!\left(\tfrac{r_i+r_{i+1}}{2}\right)&=0,\quad i=1,\dots,n-1.
\end{align*}
\end{problem}

\begin{problem}
(a) Find distinct integers $m_1,\dots,m_5$ such that $p(x)=\prod_{i=1}^5(x-m_i)$ has the minimum number of nonzero coefficients.  
(b) Let $P(x)=x^{11}+a_{10}x^{10}+\cdots+a_0$ be monic with real coefficients, $a_0\neq0$, and all real zeros. Find the least possible number of nonzero coefficients of $P(x)$.
\end{problem}

\begin{problem}
Let $P(x)$ be a polynomial of degree $n$ such that $P(x)=Q(x)P''(x)$, where $Q(x)$ is quadratic. Show that if $P(x)$ has at least two distinct roots, then it must have $n$ distinct roots.
\end{problem}

\begin{problem}
(a) Let $p(z)$ be degree $n$, all zeros $|z|=1$. Put $g(z)=p(z)/z^{n/2}$. Show all zeros of $g'(z)$ lie on $|z|=1$.  
(b) Let $f(t)=\sum_{j=1}^N a_j \sin(2\pi jt)$ with real $a_j$, $a_N\neq0$. Let $N_k$ be the number of zeros of $\tfrac{d^k}{dt^k}f(t)$ in $[0,1)$. Prove $N_0\leq N_1\leq N_2\leq\cdots$ and $\lim_{k\to\infty}N_k=2N$.
\end{problem}

\begin{problem}
For every nonconstant polynomial $p$, let $H_p=\{z\in\mathbb{C}:|p(z)|=1\}$. Prove that if $H_p=H_q$, then $p=r^m$, $q=\xi r^n$ for some $r$ polynomial, $m,n\in\mathbb{Z}^+$, and $|\xi|=1$.
\end{problem}

\begin{problem}
For each integer $m$, consider
\[
P_m(x)=x^4-(2m+4)x^2+(m-2)^2.
\]
For what values of $m$ is $P_m(x)$ the product of two nonconstant polynomials with integer coefficients?
\end{problem}

\begin{problem}
Let $k$ be fixed. The $n$th derivative of $1/(x^k-1)$ has the form $\frac{P_n(x)}{(x^k-1)^{n+1}}$. Find $P_n(1)$.
\end{problem}

\begin{problem}
Let $p$ be prime. Prove that
\[
\det\begin{pmatrix}
x & y & z \\ x^p & y^p & z^p \\ x^{p^2} & y^{p^2} & z^{p^2}
\end{pmatrix}
\]
is congruent mod $p$ to a product of polynomials of form $ax+by+cz$ with integer $a,b,c$.
\end{problem}

\begin{problem}
Let $f(z)=az^4+bz^3+cz^2+dz+e=a(z-r_1)(z-r_2)(z-r_3)(z-r_4)$, integers $a\neq0$. Show that if $r_1+r_2$ is rational and not equal to $r_3+r_4$, then $r_1r_2$ is rational.
\end{problem}

\begin{problem}
Let $P(x)=c_nx^n+\cdots+c_0$ with integer coefficients. If $P(r)=0$ for rational $r$, show that
\[
c_nr,\; c_nr^2+c_{n-1}r,\; \dots,\; c_nr^n+c_{n-1}r^{n-1}+\cdots+c_1r
\]
are integers.
\end{problem}

\begin{problem}
Let $n$ be positive. Find the number of pairs $(P,Q)$ with real coefficients such that
\[
P(x)^2+Q(x)^2=x^{2n}+1,\quad \deg P > \deg Q.
\]
\end{problem}

\begin{problem}
Let $k$ be a positive integer. Prove there exist polynomials $P_0(n),\dots,P_{k-1}(n)$ such that for any integer $n$,
\[
\left\lfloor\frac{n}{k}\right\rfloor^k = P_0(n)+P_1(n)\left\lfloor\frac{n}{k}\right\rfloor+\cdots+P_{k-1}(n)\left\lfloor\frac{n}{k}\right\rfloor^{k-1}.
\]
\end{problem}

\begin{problem}
Find the smallest $j>0$ such that for every integer-coefficient polynomial $p(x)$ and integer $k$, the derivative
\[
p^{(j)}(k)
\]
is divisible by $2016$.
\end{problem}

\begin{problem}
Let $n>0$. Show there exist positive reals $a_0,\dots,a_n$ such that for each choice of signs, the polynomial
\[
\pm a_nx^n \pm a_{n-1}x^{n-1} \pm \cdots \pm a_1x \pm a_0
\]
has $n$ distinct real roots.
\end{problem}

\begin{problem}
Determine all pairs $(P,Q)$ of complex monic polynomials such that $P(x)\mid Q(x)^2+1$ and $Q(x)\mid P(x)^2+1$.
\end{problem}

\begin{problem}
Let $p(x)$ be nonconstant real-coefficient polynomial. For $n>0$, define
\[
q_n(x)=(x+1)^np(x)+x^np(x+1).
\]
Prove only finitely many $n$ yield $q_n(x)$ with all real roots.
\end{problem}

\begin{problem}
Let $ax^3+bx^2+cx+d$ have three distinct real roots. How many real roots has
\[
4(ax^3+bx^2+cx+d)(3ax+b)=(3ax^2+2bx+c)^2?
\]
\end{problem}

\begin{problem}
Does there exist a finite set $M\subset\mathbb{R}\setminus\{0\}$ such that for every $n>0$, there exists a polynomial of degree at least $n$ with coefficients in $M$, all roots real and in $M$?
\end{problem}

\begin{problem}
Suppose $ax^2+(c-b)x+(e-d)$ has two real roots greater than $1$. Prove that $ax^4+bx^3+cx^2+dx+e$ has at least one real root.
\end{problem}

\begin{problem}
If $a,b,c\in\mathbb{C}$ such that roots of $z^3+az^2+bz+c$ all satisfy $|z|=1$, prove roots of $x^3+|a|x^2+|b|x+|c|$ all satisfy $|x|=1$.
\end{problem}

\begin{problem}
Let $P(x)=x^n+a_{n-1}x^{n-1}+\cdots+a_0$ monic with roots $x_1,\dots,x_n$. The discriminant is
\[
\Delta(P)=\prod_{1\leq i<j\leq n}(x_i-x_j)^2.
\]
Show that
\[
\Delta(x^n+ax+b)=(-1)^{\binom{n}{2}}\left(n^n b^{n-1}+(-1)^{n-1}(n-1)^{n-1}a^n\right).
\]
\end{problem}

\begin{problem}
Let $P_n(x)=(x+n)(x+n-1)\cdots(x+1)-(x-1)(x-2)\cdots(x-n)$. Show all zeros of $P_n(x)$ are purely imaginary.
\end{problem}

\begin{problem}
Let $P(x)$ have complex coefficients, all roots with real part $a$. For $|z|=1$, show every root of $R(x)=P(x-1)-zP(x)$ has real part $a+\tfrac{1}{2}$.
\end{problem}

\begin{problem}
Let $d\geq1$. Show there exists $A_d(x)$ of degree $d$ with
\[
F_d(x)=\sum_{n\geq0} n^d x^n = \frac{A_d(x)}{(1-x)^{d+1}}.
\]
For example, $A_1(x)=x$, $A_2(x)=x+x^2$, $A_3(x)=x+4x^2+x^3$. Show every root of $A_d(x)$ is real.
\end{problem}

\begin{problem}
Let $P(z)=z^n+a_{n-1}z^{n-1}+\cdots+a_0$ monic. Choose $j$ so that roots can be labeled $\alpha_1,\dots,\alpha_n$ with $|\alpha_1|,\dots,|\alpha_j|>1$, $|\alpha_{j+1}|,\dots,|\alpha_n|\leq1$. Prove
\[
\prod_{i=1}^j|\alpha_i|\leq \sqrt{|a_0|^2+\cdots+|a_{n-1}|^2+1}.
\]
\end{problem}

\begin{problem}
Let $Q(x)$ be monic degree $n$ real-coefficient. Prove
\[
\sup_{x\in[-2,2]}|Q(x)|\geq2.
\]
Optional: Prove equality only for Chebyshev polynomials.
\end{problem}

\begin{problem}
Let $P,Q$ be polynomials with real roots $r_1\leq\dots\leq r_n$, $s_1\leq\dots\leq s_{n-1}$. They are interlaced if
\[
r_1\leq s_1\leq r_2\leq s_2\leq\cdots\leq s_{n-1}\leq r_n.
\]
Prove interlacing iff $P+tQ$ has all real roots for all $t\in\mathbb{R}$.
\end{problem}

\begin{problem}
Let $P(x)$ be real-coefficient polynomial. For $t\in\mathbb{R}$, let $V(P,t)$ be sign changes in $P(t),P'(t),P''(t),\dots$. Prove: for any $a,b$, the number of roots of $P$ in $(a,b]$, counted with multiplicity, is $V(P,a)-V(P,b)$ minus a nonnegative even integer. Deduce Descartes' rule of signs.
\end{problem}

\begin{problem}
Let $P(x)$ be squarefree with real coefficients. Define $P_0=P$, $P_1=P'$, and $P_{i+2}=-\mathrm{rem}(P_i,P_{i+1})$ until constant. Prove: for any $a,b$, the number of zeros of $P$ in $(a,b]$ is $\sigma(a)-\sigma(b)$, where $\sigma(t)$ is sign changes in $P_0(t),\dots,P_r(t)$.
\end{problem}

\section*{Analysis}

\begin{problem}
Show that $\int_0^\infty \frac{\cos(ax)}{1+x^2}\,dx$ exists for $a\in\mathbb{R}$ and compute its value.
\end{problem}

\begin{problem}
Find a simple function $f(x)$ for which $x^{1/x}-1 \sim f(x)$ as $x\to\infty$.
\end{problem}

\begin{problem}
For what pairs $(a,b)$ of positive reals does the improper integral
\[
\int_b^\infty \left(\sqrt{\sqrt{x+a}}-\sqrt{x}-\sqrt{\sqrt{x}}-\sqrt{x-b}\right)\,dx
\]
converge?
\end{problem}

\begin{problem}
Let $a_n$ be the unique positive root of $x^n+x=1$. Find $f(n)$ such that $1-a_n\sim f(n)$ as $n\to\infty$.
\end{problem}

\begin{problem}
For $f:[0,1]\to\mathbb{R}$ continuous, let $I(f)=\int_0^1 x^2f(x)\,dx$, $J(f)=\int_0^1 xf(x)^2\,dx$. Find $\max(I(f)-J(f))$.
\end{problem}

\textbf{Problems.}

\begin{problem}
For $a>0$, compute $\int_0^\infty t^{-1/2} e^{-a(t+t^{-1})}\,dt$.
\end{problem}

\begin{problem}
Let $f:[0,\infty)\to\mathbb{R}$, differentiable, with $f'(x)=-3f(x)+6f(2x)$, $|f(x)|\leq e^{-\sqrt{x}}$. Define $\mu_n=\int_0^\infty x^n f(x)\,dx$.
\begin{itemize}
\item[(a)] Express $\mu_n$ in terms of $\mu_0$.
\item[(b)] Prove $\{\mu_n\cdot 3^n/n!\}$ converges, and limit is $0$ iff $\mu_0=0$.
\end{itemize}
\end{problem}

\begin{problem}
Suppose $f,g:\mathbb{R}\to\mathbb{R}$ differentiable, non-constant, with
\begin{align*}
f(x+y)&=f(x)f(y)-g(x)g(y),\\
g(x+y)&=f(x)g(y)+g(x)f(y),
\end{align*}
and $f'(0)=0$. Prove $(f(x))^2+(g(x))^2=1$.
\end{problem}

\begin{problem}
For $a,b>0$, find largest $c$ such that
\[
a^x b^{1-x}\leq \frac{a\sinh(ux)}{\sinh u}+\frac{b\sinh(u(1-x))}{\sinh u}
\]
for all $u$, $0<|u|\leq c$, and all $0<x<1$.
\end{problem}

\begin{problem}
Let $K(x,y)$ positive, continuous on $[0,1]^2$, $f,g$ positive continuous on $[0,1]$, with
\[\int_0^1 f(y)K(x,y)\,dy=g(x),\quad \int_0^1 g(y)K(x,y)\,dy=f(x).\]
Show $f(x)=g(x)$.
\end{problem}

\begin{problem}
Evaluate
\[
\int_0^\infty \left(x-\tfrac{x^3}{2}+\tfrac{x^5}{2\cdot4}-\tfrac{x^7}{2\cdot4\cdot6}+\cdots\right)\left(1+\tfrac{x^2}{2^2}+\tfrac{x^4}{2^2\cdot4^2}+\cdots\right)\,dx.
\]
\end{problem}

\begin{problem}
Let $f$ be twice differentiable, satisfying $f(x)+f''(x)=-xg(x)f'(x)$, $g(x)\geq0$. Prove $|f(x)|$ is bounded.
\end{problem}

\begin{problem}
Prove: there exists $C$ such that for any degree-1999 polynomial $p(x)$,
\[
|p(0)|\leq C\int_{-1}^1|p(x)|\,dx.
\]
\end{problem}

\begin{problem}
Find $c,L$ with
\[
\lim_{r\to\infty}\frac{r^c\int_0^{\pi/2} x^r\sin x\,dx}{\int_0^{\pi/2} x^r\cos x\,dx}=L.
\]
\end{problem}

\begin{problem}
Let $(a_i,b_i)$ be vertices of convex polygon containing origin. Prove $\exists x,y>0$ with
\[\sum (a_i,b_i)x^{a_i}y^{b_i}=(0,0).\]
\end{problem}

\begin{problem}
Show all solutions of $y''+e^xy=0$ remain bounded as $x\to\infty$.
\end{problem}

\begin{problem}
Let $f$ be real with partial derivatives on $x^2+y^2\leq1$, $|f(x,y)|\leq1$. Show $\exists(x_0,y_0)$ inside unit circle with $(\partial_x f(x_0,y_0))^2+(\partial_y f(x_0,y_0))^2\leq16$.
\end{problem}

\begin{problem}
On $[0,1]$, let $f$ continuously differentiable, $0<f'(x)\leq1$, $f(0)=0$. Prove
\[\left(\int_0^1 f(x)\,dx\right)^2\geq \int_0^1 f(x)^3\,dx,\]
and give example where equality holds.
\end{problem}

\begin{problem}
Let $P(t)$ polynomial. Prove system
\[\int_0^x P(t)\sin t\,dt=\int_0^x P(t)\cos t\,dt=0\]
has finitely many real $x$.
\end{problem}

\begin{problem}
Let $C=\{f\in C^1[0,1]: f(0)=0,f(1)=1\}$. Find largest $u$ such that
\[
u\leq \int_0^1|f'(x)-f(x)|\,dx\quad \forall f\in C.
\]
\end{problem}

\begin{problem}
Given $\sum a_n$ convergent, $a_n>0$, prove $\sum \sqrt[n]{a_1a_2\cdots a_n}$ converges.
\end{problem}

\begin{problem}
Given $f(x)+f'(x)\to0$ as $x\to\infty$, prove $f(x)\to0$, $f'(x)\to0$.
\end{problem}

\begin{problem}
Suppose $f''$ continuous, $|f(x)|\leq a$, $|f''(x)|\leq b$. Find best $c$ with $|f'(x)|\leq c$.
\end{problem}

\begin{problem}
Let $f$ have continuous third derivative, $f,f',f'',f'''\!>0$, $f'''\leq f$. Show $f'(x)<2f(x)$. 
\end{problem}

\begin{problem}
Show $\int_0^\infty \sin(x)\sin(x^2)\,dx$ converges.
\end{problem}

\begin{problem}
Fix $b\geq2$. Define $f(1)=1,f(2)=2$, $f(n)=n f(d)$ where $d=$ number of base-$b$ digits of $n$. For which $b$ does $\sum 1/f(n)$ converge?
\end{problem}

\begin{problem}
Evaluate
\[
\lim_{x\to1^-}\prod_{n=0}^\infty \left(\frac{1+x^{n+1}}{1+x^n}\right)^{x^n}.
\]
\end{problem}

\begin{problem}
Find differentiable $f:(0,\infty)\to(0,\infty)$ for which $\exists a>0$ with $f'(a/x)=x/f(x)$.
\end{problem}

\begin{problem}
Let $k>1$, $a_0>0$, $a_{n+1}=a_n+(1/k)\sqrt[k]{a_n}$. Evaluate $\lim_{n\to\infty} a_n^{k+1}/n^k$.
\end{problem}

\begin{problem}
Define $f(x)=x$ if $x\leq e$, $f(x)=xf(\ln x)$ if $x>e$. Does $\sum 1/f(n)$ converge?
\end{problem}

\begin{problem}
Find all $C^1$ $f:\mathbb{R}\to\mathbb{R}$ such that $f(q)$ rational with same denominator as $q$ for all rational $q$.
\end{problem}

\begin{problem}
Let $f,g,h$ differentiable near $0$, satisfy
\begin{align*}
f'&=2f^2gh+\tfrac{1}{gh},\quad f(0)=1,\\
g'&=fg^2h+\tfrac{4}{fh},\quad g(0)=1,\\
h'&=3fgh^2+\tfrac{1}{fg},\quad h(0)=1.
\end{align*}
Find explicit $f(x)$ near $0$.
\end{problem}

\begin{problem}
Let $f:[0,1]^2\to\mathbb{R}$ continuous, $\partial f/\partial x,\partial f/\partial y$ continuous in $(0,1)^2$. Let $a=\int_0^1 f(0,y)\,dy$, $b=\int_0^1 f(1,y)\,dy$, $c=\int_0^1 f(x,0)\,dx$, $d=\int_0^1 f(x,1)\,dx$. Must there be $(x_0,y_0)$ with $f_x(x_0,y_0)=b-a$, $f_y(x_0,y_0)=d-c$?
\end{problem}

\begin{problem}
Let $f:(1,\infty)\to\mathbb{R}$ differentiable with
\[
f'(x)=\frac{x^2-(f(x))^2}{x^2((f(x))^2+1)}.
\]
Prove $\lim_{x\to\infty} f(x)=\infty$.
\end{problem}

\begin{problem}
Find all differentiable $f:\mathbb{R}\to\mathbb{R}$ with
\[
f'(x)=\frac{f(x+n)-f(x)}{n},\quad \forall x, \forall n\in\mathbb{Z}^+.
\]
\end{problem}

\begin{problem}
Suppose $h:\mathbb{R}^2\to\mathbb{R}$ with continuous partials, $h(x,y)=a h_x(x,y)+b h_y(x,y)$. Prove: if $|h(x,y)|\leq M$, then $h\equiv0$.
\end{problem}

\begin{problem}
Let $f:[0,\infty)\to\mathbb{R}$ strictly decreasing, continuous, $\lim f(x)=0$. Show $\int_0^\infty \frac{f(x)-f(x+1)}{f(x)}\,dx$ diverges.
\end{problem}

\begin{problem}
Is there strictly increasing $f:\mathbb{R}\to\mathbb{R}$ with $f'(x)=f(f(x))$?
\end{problem}

\end{document}








